\documentclass[../../../include/open-logic-section]{subfiles}

\begin{document}

\olfileid[id]{sth}{choice}{woproblem}
\olsection{Masalah Pengurutan Baik}

Jelas bahwa cukup banyak hal bergantung pada apakah kita menerima Pengurutan
Baik. Namun, pembahasan prinsip ini cenderung berfokus pada suatu prinsip yang
ekuivalen, yaitu Aksioma Pilihan. Karena itu, sekarang kita akan mengalihkan
perhatian kepadanya (dan membuktikan ekuivalensinya).

Pada \citeyear{Cantor1883}, Cantor menyatakan dukungannya terhadap Aksioma
Pengurutan Baik dengan menyebutnya ``suatu hukum pikiran yang bagi saya
tampak mendasar, kaya akan konsekuensi, dan khususnya luar biasa karena
keberlakuannya yang umum'' (dikutip dalam \citeauthor{Potter2004}
\citeyear[hlm.~243]{Potter2004}). Namun, pada akhirnya Cantor menjadi yakin
bahwa ``Aksioma'' itu memerlukan bukti. Komunitas matematika pun
berpandangan demikian.

Masalah tersebut ``dipecahkan'' oleh Zermelo pada
\citeyear{Zermelo1904}. Untuk menjelaskan penyelesaiannya, kita memerlukan
beberapa definisi.

\begin{defn}
Suatu fungsi $f$ merupakan \emph{fungsi pilihan} jika dan hanya jika
$f(x) \in x$ bagi semua $x \in \dom{f}$. Kita mengatakan bahwa $f$
merupakan \emph{fungsi pilihan bagi $A$} jika dan hanya jika $f$ merupakan
fungsi pilihan dengan $\dom{f} = A \setminus \{\emptyset\}$.
\end{defn}

Secara intuitif, bagi setiap himpunan (tak kosong) $x \in A$, suatu fungsi
pilihan bagi $A$ \emph{memilih} satu anggota tertentu, $f(x)$, dari~$x$.
Aksioma Pilihan kemudian berbunyi:

\begin{axiom}[Pilihan]
	Setiap himpunan mempunyai suatu fungsi pilihan.
\end{axiom}

Zermelo menunjukkan bahwa Pilihan mengakibatkan pengurutan baik, dan
sebaliknya:

\begin{thm}[dalam $\ZF$]\ollabel{thmwochoice}
Pengurutan Baik dan Pilihan ekuivalen.
\end{thm}

\begin{proof}
\emph{Dari kiri ke kanan.} Misalkan $A$ suatu himpunan yang
anggota-anggotanya juga himpunan. Berdasarkan Aksioma Gabungan,
$\bigcup A$ ada; jadi, berdasarkan Pengurutan Baik, terdapat suatu $<$ yang
mengurutkan baik $\bigcup A$. Sekarang tetapkan
$f(x)=\text{anggota terkecil-$<$ dari }x$. Ini merupakan fungsi pilihan
bagi~$A$.

\emph{Dari kanan ke kiri.} Tetapkan $A$. Jika $A=\emptyset$, pengurutan
kosong mengurutkan baik~$A$; jadi, andaikan $A\neq\emptyset$. Berdasarkan
Pilihan, terdapat suatu fungsi pilihan $f$ bagi
$\Pow{A} \setminus \{\emptyset\}$. Dengan menggunakan Rekursi Transfinit,
definisikan:
\begin{align*}
	g(0) &= f(A)\\
	g(\alpha) &=
		\begin{cases}
			\text{berhenti!{}} &\text{jika }A = \funimage{g}{\alpha}\\
			f(A \setminus \funimage{g}{\alpha}) & \text{jika tidak}\\
		\end{cases}
\end{align*}
Instruksi untuk ``berhenti!'' hanya merupakan singkatan bagi definisi
yang lebih panjang: rekursi dilanjutkan selama
$A\neq\funimage{g}{\alpha}$, lalu $g$ diambil sebagai pembatasannya pada
ordinal terkecil $\alpha$ dengan $A=\funimage{g}{\alpha}$. Pada mulanya, kita
hanya dapat memastikan bahwa hal ini mendefinisikan suatu \emph{suku}
(lihat catatan sesudah
\olref[sth][spine][recursion]{transrecursionschema}); tetapi kita akan
menunjukkan bahwa kita memang memperoleh suatu fungsi.

Karena $f$ merupakan fungsi pilihan, bagi setiap $\alpha$ sebelum konstruksi
berhenti berlaku
$g(\alpha)=f(A\setminus\funimage{g}{\alpha})\in
A\setminus\funimage{g}{\alpha}$; yakni
$g(\alpha)\notin\funimage{g}{\alpha}$. Jadi, jika
$g(\alpha)=g(\beta)$, maka $g(\beta)\notin\funimage{g}{\alpha}$, yakni
$\beta\notin\alpha$, dan secara serupa $\alpha\notin\beta$. Berdasarkan
Trikotomi, $\alpha=\beta$. Jadi, $g$ bersifat !!{injective}.

Selanjutnya, perhatikan bahwa konstruksi memang berhenti; yakni terdapat
suatu ordinal terkecil $\alpha$ sedemikian sehingga
$A=\funimage{g}{\alpha}$. Andaikan sebaliknya. Karena $g$ bersifat
!!{injective}, kita akan memperoleh $\cardless{\alpha}{A}$ bagi setiap
ordinal $\alpha$, bertentangan dengan \olref[hartogs]{HartogsLemma}. Karena
itu, $\ran{g}=A$.

Dengan menggabungkan fakta-fakta ini, $g$ merupakan !!a{bijection} dari
suatu ordinal ke~$A$. Sekarang $g$ dapat digunakan untuk mengurutkan
baik~$A$.
\end{proof}

Jadi, Pengurutan Baik dan Pilihan berdiri atau runtuh bersama-sama. Namun,
pertanyaannya tetap ada: apakah keduanya berdiri atau runtuh?

\end{document}
